% Original CSC302 teaching exposition based on Macbeth, Chapter 1, Sections 1.1--1.5.
% Source: https://hrmacbeth.github.io/math2001/01_Proofs_by_Calculation.html
% Source snapshot: hrmacbeth/math2001 e660f42b13ddcb6d12b52ba036d6bd071a0cfb9b.
% Numerical variants and exercises below are newly chosen; example locators identify
% the corresponding method, not a claim that the new numbers appear in the book.
% CSC302 lecture slides adapted from Heather Macbeth, The Mechanics of Proof.
\documentclass[aspectratio=169,11pt]{beamer}
\usepackage[T1]{fontenc}
\usepackage{lmodern}
\usepackage{amsmath,amssymb,mathtools}
\usepackage{booktabs,array,listings,tikz}
\usetikzlibrary{arrows.meta,positioning,calc}
\definecolor{ProofNavy}{HTML}{18324F}
\definecolor{ProofTeal}{HTML}{007D82}
\definecolor{ProofPaper}{HTML}{F3F6F8}
\usetheme{default}
\setbeamercolor{normal text}{fg=ProofNavy,bg=white}
\setbeamercolor{structure}{fg=ProofTeal}
\setbeamercolor{frametitle}{fg=ProofNavy}
\setbeamercolor{block title}{fg=white,bg=ProofNavy}
\setbeamercolor{block body}{fg=ProofNavy,bg=ProofPaper}
\setbeamercolor{block title example}{fg=white,bg=ProofTeal}
\setbeamercolor{block body example}{fg=ProofNavy,bg=ProofPaper}
\setbeamerfont{frametitle}{size=\Large,series=\bfseries}
\setbeamertemplate{navigation symbols}{}
\setbeamertemplate{itemize item}{\small\textbullet}
\setbeamertemplate{footline}{%
  \leavevmode\hbox{%
  \begin{beamercolorbox}[wd=.86\paperwidth,ht=2.8ex,dp=1.2ex,leftskip=1.6em]{author in head/foot}%
    \scriptsize CSC302\quad\textbar\quad\insertshorttitle
  \end{beamercolorbox}%
  \begin{beamercolorbox}[wd=.14\paperwidth,ht=2.8ex,dp=1.2ex,rightskip=1.6em plus1fil]{date in head/foot}%
    \scriptsize\insertframenumber/\inserttotalframenumber
  \end{beamercolorbox}}}
\setbeamersize{text margin left=8mm,text margin right=8mm}
\lstdefinelanguage{Lean}{%
  morekeywords={example,theorem,lemma,def,by,calc,have,exact,apply,intro,intros,constructor,rcases,rfl,rw,ring,numbers,linarith,nlinarith,rel,extra,addarith,obtain,left,right,use,choose,dsimp,change,ext,fun,show,from,at,with,induction,case,decide,contradiction,push_neg,by_cases,by_contra,assumption,existsi,simp,simpa},
  sensitive=true,morecomment=[l]{--},morestring=[b]"}
\lstset{language=Lean,basicstyle=\ttfamily\small,keywordstyle=\color{ProofTeal}\bfseries,
  commentstyle=\color{gray},stringstyle=\color{ProofNavy},showstringspaces=false,
  columns=fullflexible,keepspaces=true,breaklines=true,frame=single,
  rulecolor=\color{ProofTeal!30},backgroundcolor=\color{ProofPaper},
  aboveskip=0.5em,belowskip=0.4em}
\newcommand{\N}{\mathbb{N}}
\newcommand{\Z}{\mathbb{Z}}
\newcommand{\Q}{\mathbb{Q}}
\newcommand{\R}{\mathbb{R}}
\newcommand{\coursechapter}[2]{%
  \title[Chapter #1: #2]{Chapter #1\\#2}%
  \subtitle{The Mechanics of Proof \textbar{} CSC302 lecture notes}%
  \author{Based on Heather Macbeth's textbook}%
  \date{}%
}
\newcommand{\sourcecredit}[1]{\par\vfill{\tiny Source: Heather Macbeth, \emph{The Mechanics of Proof}, #1.}}
\newcommand{\bookurl}{\url{https://hrmacbeth.github.io/math2001/}}
\hypersetup{colorlinks=true,urlcolor=ProofTeal,linkcolor=ProofNavy}

\coursechapter{1}{Proofs by calculation}
\begin{document}

\begin{frame}\titlepage\end{frame}

\begin{frame}{What you will learn}
\begin{itemize}
\item Turn hypotheses into a chain that begins and ends at the required expressions.
\item Separate algebraic identities from substitutions of known facts.
\item Control the direction of inequalities and justify sign conditions.
\item Translate individual steps into \texttt{calc}, \texttt{ring}, \texttt{rw},
      \texttt{rel}, \texttt{numbers}, \texttt{extra}, and \texttt{addarith}.
\end{itemize}
\begin{block}{A proof is an explanation}
Every link must be justified, and the links must combine to establish the stated goal.
\end{block}
\end{frame}

\begin{frame}{Start by identifying the number system}
\begin{center}
\begin{tabular}{lll}
\toprule
Domain & Examples & Lean name \\
\midrule
$\N$ & $0,1,2,\ldots$ & \texttt{Nat} \\
$\Z$ & $\ldots,-2,-1,0,1,\ldots$ & \texttt{Int} \\
$\Q$ & $-\frac32,\frac57,4$ & \texttt{Rat} \\
$\R$ & $\frac57,\sqrt2,\pi$ & \texttt{Real} \\
\bottomrule
\end{tabular}
\end{center}
\begin{itemize}
\item Here $0$ is a natural number.
\item The domain constrains the legal operations: $3-5$ in $\N$ is truncated subtraction.
\item In calculations over $\Q$ or $\R$, division by a known nonzero constant is available.
\end{itemize}
\sourcecredit{Chapter 1 introduction and Section 1.3}
\end{frame}

\section{1.1 Proving equalities}
% Method: Examples 1.1.1--1.1.2.
\begin{frame}{1.1\quad A chain proves its endpoints equal}
Suppose each of the following statements has been justified:
\[
E_0=E_1,\qquad E_1=E_2,\qquad E_2=E_3.
\]
Transitivity of equality gives $E_0=E_3$.
\begin{block}{A useful writing discipline}
Start with the left side of the goal. End with its right side.
Keep an equality sign between consecutive expressions.
\end{block}
Two common justifications:
\begin{itemize}
\item \textbf{Algebra:} the expressions agree for all values in the domain.
\item \textbf{Substitution:} a hypothesis says that one expression equals another.
\end{itemize}
\sourcecredit{Section 1.1}
\end{frame}

\begin{frame}{Worked example: reveal the expressions you know}
Let $a,b\in\Q$, with $a-b=3$ and $ab=4$. Prove $(a+b)^2=25$.
\[
\begin{aligned}
(a+b)^2
 &= (a-b)^2+4ab &&\text{algebra}\\
 &= 3^2+4\cdot4 &&\text{the two hypotheses}\\
 &=25 &&\text{arithmetic}.
\end{aligned}
\]
\begin{block}{Why does the first step hold?}
$(a-b)^2+4ab=a^2-2ab+b^2+4ab=a^2+2ab+b^2$.
\end{block}
The initial rewrite was chosen because $a-b$ and $ab$ have known values.
\sourcecredit{Section 1.1; numerical variant of Example 1.1.1}
\end{frame}

% Method and identity: Example 1.1.4. Explanation independently composed.
\begin{frame}{A calculation can avoid an unnecessary division}
Let $a,b,c,d,e,f\in\Z$, with $ad=bc$ and $cf=de$.
Then
\[
\begin{aligned}
d(af-be)&=(ad)f-dbe\\
        &=(bc)f-dbe\\
        &=b(cf)-dbe\\
        &=b(de)-dbe=0.
\end{aligned}
\]
Every step uses multiplication, subtraction, or a given equality.
\begin{block}{Check the hypotheses before cancelling}
Writing $a=bc/d$ would require $d\ne0$ and a suitable domain.
The hypotheses do not give $d\ne0$.
\end{block}
\sourcecredit{Example 1.1.4}
\end{frame}

\begin{frame}{Practice 1.1: write the chain before simplifying}
\begin{enumerate}
\item Let $r,s\in\R$, with $r+3s=2$ and $s=-1$. Prove $r=5$.
\item Let $a,b\in\Q$, with $a-b=5$ and $ab=6$. Prove $(a+b)^2=49$.
\end{enumerate}
\vspace{1em}
For each proof:
\begin{itemize}
\item State the first and last expressions before filling the middle.
\item Mark every step as algebra, substitution, or arithmetic.
\end{itemize}
\end{frame}

\begin{frame}{Practice 1.1: solution discussion}
\textbf{1.} Introduce the expression $r+3s$:
\[
\begin{aligned}
r&=(r+3s)-3s&&\text{algebra}\\
 &=2-3(-1)&&\text{substitution}\\
 &=5&&\text{arithmetic}.
\end{aligned}
\]
\textbf{2.} Use the same identity as in the worked example:
\[
(a+b)^2=(a-b)^2+4ab=5^2+4\cdot6=49.
\]
The calculation proves the requested expression directly; values of $a$ and $b$ are unnecessary.
\end{frame}

\section{1.2 Proving equalities in Lean}
% Syntax and method: Examples 1.2.1--1.2.2.
\begin{frame}[fragile]{1.2\quad Read a Lean statement as a mathematical sentence}
\begin{lstlisting}
example {a b : Rat}
    (hd : a - b = 3) (hp : a * b = 4) :
    (a + b) ^ 2 = 25 := by
\end{lstlisting}
\begin{itemize}
\item \texttt{a b : Rat}: the variables are rational numbers.
\item \texttt{hd} and \texttt{hp}: names for the hypotheses.
\item After the final colon: the proposition to prove.
\item \texttt{by}: the proof will be given by tactics.
\item Multiplication uses \texttt{*}; exponentiation uses \texttt{\^{}}.
\end{itemize}
The examples use the textbook environment: its chapter file imports
\texttt{Library.Basic} and runs \texttt{math2001\_init}.
\sourcecredit{Section 1.2 and accompanying Lean file}
\end{frame}

\begin{frame}[fragile]{The calculation, with one justification per line}
\begin{lstlisting}
example {a b : Rat}
    (hd : a - b = 3) (hp : a * b = 4) :
    (a + b) ^ 2 = 25 := by
  calc
    (a + b) ^ 2 = (a - b) ^ 2 + 4 * (a * b) := by ring
    _ = 3 ^ 2 + 4 * 4 := by rw [hd, hp]
    _ = 25 := by numbers
\end{lstlisting}
\begin{itemize}
\item \texttt{\_} repeats the preceding right side.
\item \texttt{ring} proves a polynomial identity.
\item \texttt{rw [hd, hp]} substitutes the named equalities.
\item \texttt{numbers} checks the final numerical equality.
\end{itemize}
\sourcecredit{Method of Example 1.2.1}
\end{frame}

\begin{frame}{Treat Lean feedback as a local question}
At a line of a \texttt{calc} proof, ask:
\begin{enumerate}
\item What equality must this line establish?
\item Does it hold by algebra alone, or does it use a hypothesis?
\item If it uses a hypothesis, does its left side occur in the expression?
\end{enumerate}
\begin{block}{Parentheses matter to rewriting}
$(ab)c$ and $a(bc)$ are equal, but have different syntactic structure.
Use an algebraic step to expose the exact subexpression you plan to rewrite.
\end{block}
An unfinished placeholder is still an unfinished proof. Check that all proof obligations are resolved.
\sourcecredit{Examples 1.2.2--1.2.4}
\end{frame}

\begin{frame}{Practice 1.2: select the justification}
Let $r,s\in\R$, with named hypotheses
\[
\texttt{h1}:r+3s=2,\qquad \texttt{h2}:s=-1.
\]
Fill the right column with \texttt{ring}, \texttt{rw [...]}, or \texttt{numbers}.
\begin{center}
\begin{tabular}{ll}
\toprule
Step & Justification \\
\midrule
$r=(r+3s)-3s$ & \rule{35mm}{0.3pt} \\
$(r+3s)-3s=2-3s$ & \rule{35mm}{0.3pt} \\
$2-3s=2-3(-1)$ & \rule{35mm}{0.3pt} \\
$2-3(-1)=5$ & \rule{35mm}{0.3pt} \\
\bottomrule
\end{tabular}
\end{center}
Which step would fail if you tried \texttt{ring} without using any hypothesis?
\end{frame}

\begin{frame}[fragile]{Practice 1.2: a complete Lean solution}
\begin{lstlisting}
example {r s : Real}
    (h1 : r + 3 * s = 2) (h2 : s = -1) : r = 5 := by
  calc
    r = (r + 3 * s) - 3 * s := by ring
    _ = 2 - 3 * s := by rw [h1]
    _ = 2 - 3 * (-1) := by rw [h2]
    _ = 5 := by numbers
\end{lstlisting}
\texttt{ring} cannot establish the substitution steps for arbitrary $r,s$.
The hypotheses supply their mathematical content.
\end{frame}

\section{1.3 Tips and tricks}
% Methods: Examples 1.3.2, 1.3.6--1.3.9.
\begin{frame}{1.3\quad How to find the intermediate expressions}
\begin{enumerate}
\item Write the target endpoints.
\item List useful expressions appearing in the hypotheses.
\item Rewrite the target to contain those expressions.
\item Substitute only after the matching expressions are visible.
\end{enumerate}
\begin{exampleblock}{A small example}
If $x+6=2$ in $\Z$, then
\[
x=(x+6)-6=2-6=-4.
\]
Adding and subtracting $6$ creates the expression we know.
\end{exampleblock}
\sourcecredit{Section 1.3; method of Example 1.3.2}
\end{frame}

\begin{frame}{Worked example: combine two equations}
Let $x,y\in\Q$, with $x+2y=9$ and $3x-y=13$. Prove $x=5$.
\begin{block}{Planning}
The coefficients of $y$ are $2$ and $-1$.
Adding the first expression to twice the second cancels $y$.
\end{block}
\[
\begin{aligned}
x&=\frac{(x+2y)+2(3x-y)}7&&\text{algebra}\\
 &=\frac{9+2\cdot13}7&&\text{substitution}\\
 &=5&&\text{arithmetic}.
\end{aligned}
\]
The denominator $7$ is nonzero in $\Q$.
\sourcecredit{Section 1.3; methods of Examples 1.3.6--1.3.8}
\end{frame}

\begin{frame}{Worked example: organize a polynomial around a hypothesis}
Let $w,z\in\R$ and suppose $w+2=z$. Prove
\[
w^2+5w+4=z^2+z-2.
\]
\[
\begin{aligned}
w^2+5w+4
  &=(w+2)^2+(w+2)-2&&\text{algebra}\\
  &=z^2+z-2&&\text{substitution}.
\end{aligned}
\]
\textbf{Check the rewrite:}
\[
(w+2)^2+(w+2)-2=w^2+4w+4+w+2-2=w^2+5w+4.
\]
\sourcecredit{Section 1.3; method of Example 1.3.9}
\end{frame}

\begin{frame}{Practice 1.3: create the expressions you need}
\begin{enumerate}
\item If $x+y=7$ and $2x-y=2$ in $\Q$, prove $x=3$.
\item If $a-2=b$ in $\R$, prove $a^2-a+1=b^2+3b+3$.
\end{enumerate}
\vspace{1em}
For the first problem, find a linear combination that cancels $y$.

For the second, rewrite the left side using $a-2$.
\end{frame}

\begin{frame}{Practice 1.3: solution discussion}
\textbf{1.} The $y$ coefficients already cancel:
\[
x=\frac{(x+y)+(2x-y)}3=\frac{7+2}3=3.
\]
\textbf{2.} Complete a polynomial in $a-2$:
\[
\begin{aligned}
a^2-a+1&=(a-2)^2+3(a-2)+3\\
        &=b^2+3b+3.
\end{aligned}
\]
Both proofs follow the same pattern: algebra first, substitution second.
\end{frame}

\section{1.4 Proving inequalities}
\begin{frame}{1.4\quad An inequality needs a direction and a sign condition}
For real numbers:
\begin{center}
\begin{tabular}{lll}
\toprule
Known fact & Additional condition & Valid conclusion \\
\midrule
$a\le b$ & none & $a+c\le b+c$ \\
$a\le b$ & none & $-b\le-a$ \\
$a\le b$ & $c\ge0$ & $ca\le cb$ \\
$a<b$ & $c>0$ & $ca<cb$ \\
$a\le b$ & $0\le a$ & $a^2\le b^2$ \\
\bottomrule
\end{tabular}
\end{center}
\begin{block}{Transitivity}
$A\le B<C\le D$ implies $A<D$.
A consistent chain with a strict step has a strict conclusion.
\end{block}
\sourcecredit{Section 1.4}
\end{frame}

% Numerical variant of Example 1.4.1.
\begin{frame}{Worked example: subtraction reverses the comparison}
Let $x,y\in\Z$, with $x+4\le2$ and $y+3x\ge1$. Prove $y>6$.
\[
\begin{aligned}
y&=(y+3x)-3x&&\text{algebra}\\
 &\ge1-3x&&y+3x\ge1\\
 &=13-3(x+4)&&\text{algebra}\\
 &\ge13-3\cdot2&&x+4\le2\\
 &=7>6&&\text{arithmetic}.
\end{aligned}
\]
When $x+4$ becomes larger, $13-3(x+4)$ becomes smaller.
That is why the fourth line points toward $\ge$.
\sourcecredit{Section 1.4; method of Example 1.4.1}
\end{frame}

% Numerical variant of Example 1.4.8.
\begin{frame}{Worked example: add a nonnegative square}
Let $x,y\in\R$ and suppose $x^2+y^2\le2$. Prove $(x+y)^2<5$.
\[
\begin{aligned}
(x+y)^2&\le (x+y)^2+(x-y)^2&&(x-y)^2\ge0\\
       &=2(x^2+y^2)&&\text{algebra}\\
       &\le4&&\text{the hypothesis}\\
       &<5&&\text{arithmetic}.
\end{aligned}
\]
\begin{block}{How was the extra square chosen?}
Its expansion cancels the $2xy$ term and produces the known expression $x^2+y^2$.
\end{block}
\sourcecredit{Section 1.4; method of Example 1.4.8}
\end{frame}

\begin{frame}{Valid individual comparisons can form an invalid proof}
If $n\in\Z$ and $n\ge6$, all these statements are true:
\[
n^2\ge36>28\le3n+10.
\]
They do \textbf{not} combine to prove $n^2>3n+10$: the direction changes.
For comparison, $4>3\le100$ does not imply $4>100$.
\begin{block}{A correct chain}
Since $n\ge6>0$,
\[
n^2=n\cdot n\ge6n=3n+3n\ge3n+18>3n+10.
\]
\end{block}
Keep some dependence on $n$ so that the final expression remains comparable.
\sourcecredit{Section 1.4; method of Example 1.4.6}
\end{frame}

\begin{frame}[fragile]{Lean follows the same inequality reasoning}
\begin{lstlisting}
example {x y : Real} (h : x ^ 2 + y ^ 2 <= 2) :
    (x + y) ^ 2 < 5 := by
  calc
    (x + y) ^ 2 <= (x + y) ^ 2 + (x - y) ^ 2 := by extra
    _ = 2 * (x ^ 2 + y ^ 2) := by ring
    _ <= 2 * 2 := by rel [h]
    _ < 5 := by numbers
\end{lstlisting}
\texttt{extra}: elementary order facts such as adding a square.\par
\texttt{rel [h]}: transport a known comparison through compatible operations.\par
\texttt{rw} is for equality substitution; a hypothesis with $\le$ needs order reasoning.
\sourcecredit{Examples 1.4.7--1.4.8}
\end{frame}

\begin{frame}{Practice 1.4: retain useful terms}
\begin{enumerate}
\item Let $x\in\R$, with $x\ge4$. Prove $x^2-2x\ge8$.
\item Let $a,b\in\R$. Prove $a^2+b^2\ge2ab$.
\end{enumerate}
\vspace{1em}
For each inequality, identify the nonnegative factor or square that makes the key step valid.
\end{frame}

\begin{frame}{Practice 1.4: solution discussion}
\textbf{1.} We know $x\ge4>0$, so multiplying by $x$ preserves $\ge$:
\[
x^2-2x\ge4x-2x=2x\ge8.
\]
\textbf{2.} Use the square of the difference:
\[
a^2+b^2=(a-b)^2+2ab\ge2ab.
\]
The second result has no sign assumptions on $a$ or $b$;
the nonnegativity comes from $(a-b)^2$.
\end{frame}

\section{1.5 A shortcut}
\begin{frame}[fragile]{1.5\quad Abbreviate additive rearrangements}
The book permits a short justification when a deduction needs only
addition or subtraction of terms on the two sides.
\begin{lstlisting}
example {x : Int} (h : x + 6 = 2) : x = -4 := by
  addarith [h]

example {u v : Real} (h : u <= 9 - v) : u + v < 10 := by
  addarith [h]
\end{lstlisting}
On paper: ``From $u\le9-v$, we get $u+v\le9<10$.''
\begin{block}{The course convention}
For deductions requiring multiplication or division, keep the relevant calculation visible.
\end{block}
\sourcecredit{Section 1.5}
\end{frame}

\begin{frame}{Practice 1.5: which steps are merely additive?}
Assume the variables are real numbers.
\begin{enumerate}
\item $a-2b=3\quad\Longrightarrow\quad a=2b+3$.
\item $4w+1=9\quad\Longrightarrow\quad w=2$.
\item $m\le7-n\quad\Longrightarrow\quad m+n<8$.
\end{enumerate}
\vspace{1em}
Explain each deduction. For the one that requires division, write a full chain.
\end{frame}

\begin{frame}{Practice 1.5: solution discussion}
\textbf{1.} Add $2b$ to both sides: an additive rearrangement.
\[
a=2b+3.
\]
\textbf{2.} Division by $4$ is essential:
\[
w=\frac{(4w+1)-1}{4}=\frac{9-1}{4}=2.
\]
\textbf{3.} Add $n$ to both sides and compare constants:
\[
m+n\le7<8.
\]
The presence of $2b$ in the first statement does not mean that the deduction multiplies both sides.
\end{frame}

\begin{frame}{Chapter 1: a checklist for calculation proofs}
\begin{enumerate}
\item Identify the domain, hypotheses, and target endpoints.
\item Expose known expressions through algebraic identities.
\item Give a reason for each equality or inequality.
\item Check nonzero denominators and signs of multipliers.
\item Check that the relations combine in one direction.
\item Use Lean feedback to inspect individual steps.
\end{enumerate}
\begin{block}{Next chapter}
A proof can branch, build intermediate facts, and introduce witnesses.
Calculations will become components of larger arguments.
\end{block}
\end{frame}

\begin{frame}{Reading and source references}
Heather Macbeth, \emph{The Mechanics of Proof}, Chapter 1:
\begin{itemize}
\item 1.1 Proving equalities
\item 1.2 Proving equalities in Lean
\item 1.3 Tips and tricks
\item 1.4 Proving inequalities
\item 1.5 A shortcut
\end{itemize}
\small
\url{https://hrmacbeth.github.io/math2001/01_Proofs_by_Calculation.html}
\par\medskip
Companion code: \url{https://github.com/hrmacbeth/math2001}
\par\medskip
These slides provide an independently written classroom presentation and new numerical practice.
Use the book for additional examples and exercises.
\end{frame}
\end{document}
